\scanpage{72}\mspage{59}
% archive sheet number: 71

Contrairement à ce que prétend Anderson,
il n'est pas vrai en général que ces
deux foncteurs soient adjoints, sans hypothèses
supplémentaires. Voici un
\unsure{contreexemple|contre-exemple} : on prend $C=C'$, $u_!=u^*=\mathrm{id}_C$,
$W'=\mathrm{Is}_{\deleted{\unsure{e|C}}}(C)$, $W=\mathrm{Fl}(C)$, donc \deleted{\unsure{le}} $u_W^{\unsure{'|1|*}}$ est
le foncteur canonique

\diagram{\begin{tikzcd}[ampersand replacement=\&,column sep=large]
\underbrace{\Pi C}_{\substack{\text{groupoïde}\\\text{fondamental}}} \& C\, . \arrow[l]
\end{tikzcd}}

Supposons que $C$ soit 1-connexe, \deleted{\unsure{donc}} i.e
$\deleted{\unsure{C}}\Pi C$ soit équiv. à $e$. Soit $e\in\mathrm{Ob}\,C$
tel que \deleted{le \unsure{groupe} enveloppant du
monoïde} $\mathrm{End}_C(e)$ soit le groupe unité,
\deleted{p. ex. tel que $\mathrm{End}_C(e)=\{\mathrm{id}_e\}$,} \deleted{Alors} et
prenons $C_0=\{e\}$. Alors l'hypothèse
sur $C_0$ est satisfaite. Mais si $u_!^W$
était adjoint à gauche \inserted{de $u_W^*$}, $e$ serait objet
initial de $C$ ; \deleted{\unsure{donc} des objets rigides}
par exemple si $C$ a un objet final $e$
(ce qui implique que $C$ est 1-connexe et
$e$ rigide), cet objet final devrait
être initial si Anderson avait raison avec
``it is clear that \ldots''.
